\documentclass[11pt,a4paper]{article}
\usepackage[T1]{fontenc}
\usepackage[utf8]{inputenc}
\usepackage{lmodern,amsmath,amssymb}
\usepackage[margin=25mm]{geometry}
\usepackage{longtable,booktabs,array,calc}
\usepackage[hidelinks]{hyperref}
\hypersetup{pdftitle={Atlas of halving: what one refinement step does, and what emerges, in each dimension},pdfauthor={navstokgap research working notes}}
\newcommand{\tightlist}{\setlength{\itemsep}{0pt}\setlength{\parskip}{0pt}}
\setlength{\emergencystretch}{3em}
\setcounter{secnumdepth}{0}
\begin{document}
\hypertarget{atlas-of-halving-what-one-refinement-step-does-and-what-emerges-in-each-dimension}{%
\section{Atlas of halving: what one refinement step does, and what
emerges, in each
dimension}\label{atlas-of-halving-what-one-refinement-step-does-and-what-emerges-in-each-dimension}}

\textbf{Purpose, 2026-09-27 (user direction: a collective understanding
of the halving of the lattice and what emerges in each case and
dimension, as the mid-term goal).} This note is the shared map. Each
cell says what one halving does, what it leaves behind, and what
survives the limit, with the status of the statement and the note that
holds it. It adds no new theorem; it fixes one vocabulary so that the
agents working here (Claude, GPT-6 Astra, Fable reviewers) and the user
fill the same cells. Open cells are listed in §5.

Status labels: \textbf{proved} (theorem in a repository note, reviewed
or not as marked), \textbf{exact} (closed-form identity),
\textbf{formal} (derivation assuming an unproved estimate),
\textbf{known} (established literature, cited), \textbf{open}.

\hypertarget{the-operation}{%
\subsection{1. The operation}\label{the-operation}}

\textbf{Halving one direction} \(a_1\mapsto a_1/2\) of a lattice with
plaquette heat times
\(t_{\mu\nu}=\lambda_Da_\mu a_\nu/\prod_{\rho\ne\mu,\nu}a_\rho\)
(\href{series-parallel-gauge-refinement.md}{series/parallel note}, eq.
(1), small-field regime) splits exactly into two moves.

\begin{itemize}
\tightlist
\item
  \textbf{Series move.} A face containing the refined direction is cut
  in two; its heat time halves. Integrating the cutting edge is a
  convolution of the two half-weights: character coefficients multiply.
  For heat-kernel weights it closes exactly in every dimension.
\item
  \textbf{Parallel move.} A face transverse to the refined direction
  gets a new copy in the mid-plane, both at doubled heat time.
  Integrating the mid-plane leaves the factor \(\Psi\): a
  \((D-1)\)-dimensional gauge theory whose edges carry Brownian-bridge
  laws fixed by the coarse field. Pointwise products of weights add
  their logarithms; the heat kernel is closed under this only up to
  vortex terms (\href{zero-spacing-any-action.md}{zero-spacing note},
  Proposition 1).
\item
  \textbf{Count.} A halving produces parallel moves in \(\binom{D-1}2\)
  transverse planes. An isotropic step is \(D\) directional halvings;
  the isotropic heat time scales as \(t\propto\lambda_Da^{4-D}\).
\item
  \textbf{Time-only halving} (Euclidean time at fixed spatial lattice)
  sends the transverse heat times to infinity; the Hamiltonian limit
  needs exponential (Wilson-type) magnetic weights (zero-spacing note,
  §1).
\end{itemize}

\textbf{Newton's halving} is the one-dimensional case: inserting a time
\(t_{2.6}\) in a cell is a series move
(\href{refinement-composition-and-limit.md}{refinement note}, §2). A
physical record at the inserted time is a new variable attached to the
body, the analogue of a parallel insertion. Forgetting a Gaussian record
is an exact parallel move with defect proportional to \(\hbar^2\), the
identity for \(\hbar=0\) (\href{newton-record-parallel-move.md}{record
as parallel move}).

\hypertarget{b.-cuts-at-any-position-user-remark-2026-09-27}{%
\subsubsection{1b. Cuts at any position (user remark,
2026-09-27)}\label{b.-cuts-at-any-position-user-remark-2026-09-27}}

Halving is the easiest cut; any cut at fraction \(s\in(0,1)\) of a cell
performs the same two moves, with \(s\) as a parameter.

\begin{itemize}
\tightlist
\item
  \textbf{Series move.} The cut face splits into heat times \(st\) and
  \((1-s)t\); the convolution returns \(t\) exactly, for heat-kernel
  weights, in every dimension.
\item
  \textbf{Parallel move.} In \(D=3\) the transverse copies get heat
  times \(t/s\) and \(t/(1-s)\); at the quadratic level precisions add,
  \(s/t+(1-s)/t=1/t\). The inserted edges carry the bridge law observed
  at fraction \(s\), variance \(s(1-s)t_e\) in place of \(t_e/4\).
  Halving maximizes the bridge variance, so each non-symmetric cut is
  gentler than a halving; the \(s\)-dependence of the defect bounds
  (Prop. 2, Theorem 5) is expected through \(4s(1-s)\) and has not been
  derived.
\item
  \textbf{Newton.} The cell action \(K_\tau=F^2\tau^3/(24m)\)
  (\href{galileo-two-path-interference.md}{two-path note}) obeys
  \(\tau^3=(s\tau)^3+((1-s)\tau)^3+3s(1-s)\tau^3\), so one cut removes
  exactly \(3s(1-s)K_\tau\), at most \(\frac34K_\tau\) (halving). Along
  any sequence of cuts whose mesh tends to zero the removed amounts sum
  to \(K_\tau\), whatever the positions and order (the remainder is
  \(\sum_i C\tau_i^3\le C\tau\max_i\tau_i^2\)). Proved, with the
  Cameron--Martin reading (each share is the energy of the Schauder hat
  the cut inserts), in the \href{cut-measure-newton.md}{cut-measure
  note}. The Galileo action is therefore an additive measure on the cut
  process: the refinement limit is cut-independent, and every cut spends
  a definite share of the action that Newton sends to zero. This is the
  Lévy--Ciesielski structure of the Brownian bridge (variance \(s(1-s)\)
  at each dyadic or non-dyadic insertion) seen in the action.
\item
  \textbf{The rod} (user remark, 2026-09-27). The dialecticians' stick
  of \emph{Zhuangzi} 33, 一尺之捶，日取其半，萬世不竭 (``a one-foot
  stick, each day take half, in ten thousand generations it is not
  exhausted'';
  \href{../docs/classics/Zhuangzi_33_Tianxia_zh_wikisource.md}{text},
  attributed to the 辯者, the debaters, and not to Hui Shi personally),
  is Zeno's dichotomy as a cut schedule: every day one cut, at
  \(s=\frac12\) of the remaining piece only. The mesh stays at
  \(\tau/2\), so this schedule is no refinement, and the action identity
  above prices it exactly: day \(k\) removes
  \(\frac34K_\tau8^{-(k-1)}\), the total spent is
  \(\frac34\cdot\frac87K_\tau=\frac67K_\tau\), and \(\frac17K_\tau\)
  stays in the cut-off pieces. The stick is inexhaustible and so is the
  action: \(\frac17\) of it is never reached by this schedule. The
  remainder converges to a point, and the Mohist Canon B names that
  point: 非半弗𣃈，則不動，說在端 (``without halving there is no cutting
  and no moving; the explanation lies in the 端, the point''), with the
  explanation 前則中無爲半，猶端也。前後取，則端中也
  (\href{../docs/classics/Mozi_Canons_B_Explanations_JingShuoXia_zh_wikisource.md}{text};
  working gloss: ``taking from the front, the middle is never halved,
  and one ends as at a point; taking from front and back, the point is
  in the middle''). The two Mohist protocols, one-ended and two-ended,
  both converge to a point and neither refines the whole stick.
  Refinement, which spends all of \(K_\tau\), needs every piece cut
  again. Resemblance with Zeno is recorded as convergence; transmission
  to China would need evidence this note does not have.
\item
  \textbf{The frame.} A cut at an arbitrary, even irrational, position
  is the geometers' cut of Book I's closing scholium. The questions for
  each cell become: does the limit depend on the cut sequence (for the
  free field and for Newton, no), and does the universal part (the
  \(\log2\) of cell 4, which for general cuts should read \(\log(1/s)\)
  summed over the step) survive when \(s\) varies (conjecture: yes, as
  the scheme- independent one-loop coefficient).
\end{itemize}

\hypertarget{the-atlas-by-dimension}{%
\subsection{2. The atlas by dimension}\label{the-atlas-by-dimension}}

\begin{longtable}[]{@{}
  >{\raggedright\arraybackslash}p{(\columnwidth - 10\tabcolsep) * \real{0.1667}}
  >{\raggedright\arraybackslash}p{(\columnwidth - 10\tabcolsep) * \real{0.1667}}
  >{\raggedright\arraybackslash}p{(\columnwidth - 10\tabcolsep) * \real{0.1667}}
  >{\raggedright\arraybackslash}p{(\columnwidth - 10\tabcolsep) * \real{0.1667}}
  >{\raggedright\arraybackslash}p{(\columnwidth - 10\tabcolsep) * \real{0.1667}}
  >{\raggedright\arraybackslash}p{(\columnwidth - 10\tabcolsep) * \real{0.1667}}@{}}
\toprule\noalign{}
\begin{minipage}[b]{\linewidth}\raggedright
\(1+n\)
\end{minipage} & \begin{minipage}[b]{\linewidth}\raggedright
Series move
\end{minipage} & \begin{minipage}[b]{\linewidth}\raggedright
Parallel planes per halving
\end{minipage} & \begin{minipage}[b]{\linewidth}\raggedright
What one step leaves
\end{minipage} & \begin{minipage}[b]{\linewidth}\raggedright
What survives the limit
\end{minipage} & \begin{minipage}[b]{\linewidth}\raggedright
Status and notes
\end{minipage} \\
\midrule\noalign{}
\endhead
\bottomrule\noalign{}
\endlastfoot
\(0+0\) & none (a single weight) & none & the weight \(k_t(U)\) itself;
a single integral with an \(e^{-c/\hbar}\) structure & a large-\(N\)
transition (Gross--Witten) & known; \href{dimension-ladder.md}{dimension
ladder} \\
\(1+0\) Newton & exact after one cubic counterterm \(-F^2h^3/(24M)\);
the defect \(-F^2uvh/(8M)\) is a coboundary & none & a scalar &
unobserved: the exact propagator at every \(\hbar\); recorded: a floor
only if joint determinacy fails & proved;
\href{refinement-composition-and-limit.md}{refinement note},
\href{principia-fifth-postulate.md}{fifth postulate} \\
\(1+1\) & exact (heat-kernel convolution) & 0 & nothing for heat
kernels; for any action, \(a^{-2}(1-\hat c_R)\) & a
conjugation-invariant Lévy exponent \(\psi(R)\); Yang--Mills iff
Lindeberg; circle spectrum \(\hbar cL\psi(R)\); no local particle &
proved (Theorem 2 of the zero-spacing note); known (Lévy) \\
\(1+2\) & exact & 1 & free: an irrelevant quadratic form of relative
size \(a^2K^2/16+(Ga)^2/8\); \(U(1)\): vortices of density
\(e^{-\pi^2/(8t)}\); \(SU(N)\): coupling shifts \(O(t)\), curvature term
explicit & free photon for \(U(1)\) at fixed \(\lambda_3\) (every
lattice gap dies); conjectured gap \(C_3\hbar c\lambda_3\) for \(SU(N)\)
& proved for free and \(U(1)\); formal for \(SU(N)\) (Proposition 7);
see §3 \\
\(1+3\) & exact & 3 & \(t=g^2\) fixed: coupling shifts that recur each
step; large fields suppressed only as \((a\Lambda)^{2b_0c}\) & running
coupling, \(\Lambda=\mu e^{-1/(2b_0g^2)}\); small instantons die
(\(11N/3>4\)); dislocations need \(c>2/b_0\) (\(\tfrac{6}{11}\) of an
instanton for \(SU(2)\)); \(U(1)\): Coulomb phase & known (running, 6/11
criterion, Driver); the measure-level route fails for \(U(1)\)
(proved) \\
\(1+n\), \(n\ge4\) & exact & \(\binom n2\) & \(t\) grows as \(a\to0\) &
no small-field regime; a phase transition, no established continuum
limit & known (Creutz for \(SU(2)\) in 5D) \\
\end{longtable}

The same heat time \(t=\hbar g_{\rm cl}^2a^{4-D}\) controls the column
``what one step leaves''. For \(D<4\) refining and taking \(\hbar\to0\)
push it the same way; in mechanics refining makes each cell more quantum
(\href{dimension-ladder.md}{dimension ladder}, §2).

\hypertarget{the-atlas-for-12-by-group}{%
\subsection{\texorpdfstring{3. The atlas for \(1+2\) by
group}{3. The atlas for 1+2 by group}}\label{the-atlas-for-12-by-group}}

\begin{longtable}[]{@{}
  >{\raggedright\arraybackslash}p{(\columnwidth - 8\tabcolsep) * \real{0.2000}}
  >{\raggedright\arraybackslash}p{(\columnwidth - 8\tabcolsep) * \real{0.2000}}
  >{\raggedright\arraybackslash}p{(\columnwidth - 8\tabcolsep) * \real{0.2000}}
  >{\raggedright\arraybackslash}p{(\columnwidth - 8\tabcolsep) * \real{0.2000}}
  >{\raggedright\arraybackslash}p{(\columnwidth - 8\tabcolsep) * \real{0.2000}}@{}}
\toprule\noalign{}
\begin{minipage}[b]{\linewidth}\raggedright
Group
\end{minipage} & \begin{minipage}[b]{\linewidth}\raggedright
Isolated cube
\end{minipage} & \begin{minipage}[b]{\linewidth}\raggedright
Full mid-plane
\end{minipage} & \begin{minipage}[b]{\linewidth}\raggedright
One step, measure level
\end{minipage} & \begin{minipage}[b]{\linewidth}\raggedright
Iteration and limit
\end{minipage} \\
\midrule\noalign{}
\endhead
\bottomrule\noalign{}
\endlastfoot
free (\(\mathbb R^n\)) & exact & exact defect (Proposition 2 of the
series/parallel note) & Gaussian, exact & blocked actions converge to a
Gaussian fixed point (known, Bell--Wilson; for this blocking
asserted) \\
\(U(1)\) & exact (Proposition 3) & exact charge form, vortex bound
(Proposition 4, Corollary) & one step equals the free step up to density
\(e^{-\pi^2/(8t)}\) (Theorem 5, refereed); whole measure within TV
\(2(L/a)^3e^{-\pi^2/(6\lambda_3a)}\) of the monopole-free part
(\href{villain-monopole-refinement.md}{monopole note}) & free photon
(Gross 1983, known) \\
\(SU(2)\) & midpoint characters exact for every spin; spin-\(\frac12\)
and spin-1 cube sectors exact, with matrix image bounds for spin 1
(\href{su2-midpoint-exact.md}{closed form}, Theorems 1--4) &
\href{su2-midplane-order-t.md}{formal order-\(t\) expansion}: cut shifts
positive, transverse negative; other local operators have dimension at
least six; explicit torus winding term & formal; needs normalized
estimate (18) with winding control & open \\
\(SU(3)\) & curvature term explicit, softening \(t_j|X_j|^2/512\)
(group-general form of Proposition 6) & open & formal & open; the gap
\(C_3\hbar c\lambda_3\) is the infrared obligation \\
\end{longtable}

\hypertarget{what-emerges-read-across-the-atlas}{%
\subsection{4. What emerges, read across the
atlas}\label{what-emerges-read-across-the-atlas}}

\begin{itemize}
\tightlist
\item
  \textbf{Series moves never produce anything new.} In every dimension
  they close exactly for heat-kernel weights, and in mechanics after one
  scalar counterterm. The one-dimensional Newton problem and the
  two-dimensional gauge theory are all series.
\item
  \textbf{Parallel moves carry everything that renormalizes.} Their
  number, \(\binom{D-1}2\), is zero exactly where the theory is exactly
  soluble.
\item
  \textbf{The exponential ladder.} Each compact correction left by a
  step is \(e^{-c/t}\) (vortices, monopoles, dislocations). In \(D<4\)
  it is summable per physical volume at fixed coupling, so compact
  effects die; in \(D=4\) the running turns it into a power of \(a\), so
  survival is a threshold condition; in \(D>4\) there is no small \(t\).
\item
  \textbf{What survives is trajectory-dependent.} Compact \(U(1)\) in
  \(1+2\) keeps a gap along \(\lambda_3a\simeq c_0/(2\log(1/a))\),
  \(c_0\approx4.99\) the monopole exponent, where the monopole density
  per physical volume diverges, and loses it at fixed \(\lambda_3\)
  (Göpfert--Mack versus Gross).
\item
  \textbf{Constants that emerge.} Mechanics: one scalar counterterm per
  cell, and one action constant if joint determinacy is denied (Theorems
  B and B\('\) of the fifth-postulate note). \(1+1\): the Lévy exponent
  (one coupling for Yang--Mills). \(1+2\): one coupling, shifted
  summably. \(1+3\): one scale \(\Lambda\) by dimensional transmutation.
\end{itemize}

\hypertarget{open-cells}{%
\subsection{5. Open cells}\label{open-cells}}

\begin{enumerate}
\def\labelenumi{\arabic{enumi}.}
\tightlist
\item
  \(SU(2)\) cube: \emph{filled through \(J=1\)}
  (\href{su2-midpoint-exact.md}{Theorem 4}, exact diagonal entries,
  image bounds and cube sector); \(J>1\) entries and control of the full
  spin sum remain open.
\item
  \(SU(2)\) full mid-plane in \(1+2\):
  \href{su2-midplane-order-t.md}{formal order-\(t\) calculation filled},
  refereed (all items accepted): cut shifts positive, transverse
  negative, remaining bulk terms of dimension six; the fixed-torus
  winding term is removed by the clause \(\min_jN_j\ge c_0/t\); open:
  the uniform estimate (18).
\item
  Stability of Hypothesis P(\(\alpha\)) under iteration. \emph{For
  \(U(1)\) in \(1+2\), reduced to exact Gaussian blocking} (Corollary
  2\('\) of the \href{villain-monopole-refinement.md}{monopole note});
  open for \(SU(N)\).
\item
  \(1+3\): the one-loop coefficient \(-2b_0\log2\) per isotropic step
  from the three parallel insertions of each halving.
\item
  Time-only halving. \emph{Structural reading, 2026-09-27, of known
  results.} At fixed spatial lattice the electric faces make the series
  moves: heat-kernel electric weights are a semigroup in the time step,
  so they compose exactly. The magnetic faces make the parallel moves:
  with exponential weights \(e^{-a_0V}\) they close exactly in their own
  family, since products of such weights add the exponents. The whole
  defect of a time halving is the non-commutation of the two families,
  and it vanishes as \(a_0\to0\) by the Trotter product formula
  (\href{https://doi.org/10.1090/S0002-9939-1959-0108732-6}{Trotter
  1959}; \href{https://doi.org/10.1016/0022-1236(68)90020-7}{Chernoff
  1968}; metadata; strong convergence for the Laplacian on \(G^E\) plus
  a bounded magnetic potential). This is why the Hamiltonian limit
  (\href{https://doi.org/10.1103/PhysRevD.11.395}{Kogut and Susskind
  1975}; transfer matrix:
  \href{https://doi.org/10.1103/PhysRevD.15.1128}{Creutz 1977};
  metadata) exists at every fixed spatial lattice, while spatial halving
  leaves the non-closing factor \(\Psi\). Open: an operator-norm rate,
  which needs domain estimates for the commutator of the Laplacian with
  \(V\).
\item
  \emph{Filled 2026-09-27 for Gaussian records}
  (\href{newton-record-parallel-move.md}{record as parallel move}):
  forgetting a record convolves momentum with variance
  \(\hbar^2/(4\sigma^2)\); precisions add in parallel; refinement
  converges iff \(\sum\sigma_j^{-2}<\infty\); identity for \(\hbar=0\).
\item
  The infrared: from ultraviolet control in \(1+2\) to \(C_3>0\).
\end{enumerate}

\hypertarget{consequence-for-state}{%
\subsection{6. Consequence for STATE}\label{consequence-for-state}}

STATE carries the mid-term goal and points here. Each result that fills
a cell updates the cell in this note and nothing else; the notes that
prove results remain the record.
\end{document}
